\documentclass[11pt,a4paper]{article}
\usepackage{notes}
\usepackage{fancyhdr}

\pagestyle{fancy}
\fancyhf{}
\lhead{Geophysics 3A: Potential Fields \& Geothermics}
\rhead{Week 6 Magnetic Susceptibility}
\cfoot{\thepage}

\title{Week 6 --- Demo: Magnetic Susceptibility of Rocks}
\author{Geophysics 3A: Potential Fields \& Geothermics}
\date{\today}

\begin{document}
\maketitle

\section*{Purpose}

This activity links magnetic susceptibility to mineralogy and geological processes such as alteration, metamorphism, and weathering.

\section*{Key Reminder}

Magnetic susceptibility ($\chi$) describes how strongly a material becomes magnetised:
\[
\mathbf{M} = \chi \mathbf{H}
\]

\section*{Samples and Observations}

\begin{table}
    \centering
    \begin{tabular}{|c|p{2.5cm}|p{2.5cm}|c|p{3cm}|}
        \toprule
        \textbf{Sample} &
        \textbf{Rock type} &
        \textbf{Geological context} &
        \textbf{Susceptibility (SI)} &
        \textbf{Notes} \\
        \midrule
        1 & & & & \\
        2 & & & & \\
        3 & & & & \\
        4 & & & & \\
        5 & & & & \\
        6 & & & & \\
        7 & & & & \\
        8 & & & & \\
        \bottomrule
    \end{tabular}
\end{table}

\section*{Paired Comparisons}

\subsection*{Peridotite vs Serpentinite}
How does serpentinisation lead to high magnetic susceptibility?

\vspace{3cm}

\subsection*{Fresh Gabbro vs Altered Basalt}
Why can alteration dramatically reduce susceptibility in mafic rocks?

\vspace{3cm}

\subsection*{Granite or Sedimentary Protolith vs Metamorphosed Equivalent}
Which minerals control susceptibility in felsic or sedimentary rocks?

\vspace{3cm}

\section*{Synthesis}

\begin{enumerate}
  \item Why is susceptibility not a simple proxy for bulk composition?
  \item Which geological processes increase susceptibility?
  \item Which decrease it?
\end{enumerate}

\section*{Reflection}

Explain why petrology is essential for interpreting magnetic anomalies.

\end{document}